\section{Literal receiver and arithmetic domain}

Let \(p\equiv1\pmod {12}\) be prime and let
\[
 \frac4p=\frac1a+\frac1Y+\frac1Z,
 \qquad a,Y,Z\in\mathbb Z_{>0},\qquad Y<Z,
\]
be an existing Erd\H{o}s--Straus witness, written in the complete first-half
coordinates of the preceding global-sign theorem. Thus \(a\) is the unique
smallest denominator,
\[
 \frac p4<a<\frac p2,\qquad R=4a-p,\qquad 3\le R<p,\qquad w=\frac pR.
\]
Every witness lies in exactly one of the following two coordinate channels,
up to the retained exchange of the final two denominators:
\begin{align}
 E:\quad
 (Y,Z)&=\left(\frac{pa+a^2/u}{R},\frac{pa+p^2u}{R}\right),
 &u\mid a^2,\quad R\mid4u+1,\label{eq:Echart}\tag{E}\\
 M:\quad
 (Y,Z)&=(pb,pc),
 &b\ge2,\quad c\ge b+1,\quad
 a(4bc-b-c)=pbc.\label{eq:Mchart}\tag{M}
\end{align}
These are the original integral coordinates, not continuously enlarged model
domains.

Retain the ordered literal roots
\[
 \ell=(p,a,Y,Z),\qquad S=p+a+Y+Z,
\]
and define
\[
 H(T)=-S^{-1}\prod_{j=0}^3(T-\ell_j)
     =AT^4+T^3+BT^2+CT+D,\qquad A=-S^{-1}.
\]
Put \(d_j=H'(\ell_j)\). For each label choose and retain one square root
\(\xi_j\in\mathbb C^\times\) satisfying
\(\xi_j^2d_j=1\). The literal odd receiving map
\(O:\mathbb C^4_{\mathrm{labels}}\to\mathbb C^4_q\) has columns
\begin{equation}\label{eq:receiver}
 o_j=\xi_j
 \begin{pmatrix}
 1\\-i d_j\\A d_j^2+2\ell_jd_j\\
 i(7\ell_j^2d_j-13d_j^2)
 \end{pmatrix}.
\end{equation}
All norms below are the operator norms induced by the displayed standard
Hermitian coordinates. A change of metric is treated separately in
Section~\ref{sec:transport}.

Let \(e_r\) be the elementary symmetric functions of the four literal roots
and set
\begin{align}
 \mathfrak N={}&60e_3S^5-20e_2^2S^4+87e_2^3S^2-282e_2e_3S^3-51e_4S^4\notag\\
 &-28e_2^4+22e_2^2e_3S+431e_3^2S^2+204e_2e_4S^2\notag\\
 &+224e_2^2e_4-168e_2e_3^2-856e_3e_4S-448e_4^2.\label{eq:N}
\end{align}
The preceding exact reduction of the four rows in \eqref{eq:receiver} gives
\begin{equation}\label{eq:detO}
 \det O=-\frac{4\epsilon\mathfrak N}{S^3},
 \qquad
 \epsilon=A^2\prod_{i<j}(\ell_j-\ell_i)\prod_j\xi_j\in\{1,-1\}.
\end{equation}
The sign records both root order and all four square-root choices.

Define the exact constants
\begin{equation}\label{eq:constants}
 \mathfrak c_0=\frac{81920}{483^6},
 \qquad
 \mathfrak c=\frac{80}{6279^4}
 =\frac{81920}{4\,52^4\,483^4}.
\end{equation}

\begin{theorem}[Sharp prime growth of the literal inverse]\label{thm:main}
Every positive integral witness in the domain above satisfies
\begin{align}
 \|O^{-1}\|_2
 &<\frac{360}{\mathfrak c}\,p\sqrt w
 \le\frac{120\sqrt3}{\mathfrak c}\,p^{3/2},\label{eq:inverse-main}\\
 \|\wedge^2O^{-1}\|_2
 &<\frac{180}{\mathfrak c},\label{eq:wedge2-main}\\
 \|\wedge^3O^{-1}\|_2
 &<\frac{41}{4\mathfrak c\,p^2},\label{eq:wedge3-main}\\
 |\det O^{-1}|
 &<\frac1{4\mathfrak c\,p^2S^3}
 <\frac1{4\mathfrak c\,p^{13/2}}.\label{eq:detinv-main}
\end{align}
The exponents in \eqref{eq:inverse-main}--\eqref{eq:wedge3-main} are optimal
over actual prime witnesses.

If \(s_1(O)\ge s_2(O)\ge s_3(O)\ge s_4(O)>0\), then
\[
 s_3(O)>\sqrt{\mathfrak c/180}.
\]
Consequently the singular-value decomposition gives
\begin{equation}\label{eq:rankone}
 O^{-1}=u\varphi+B,
 \qquad \operatorname{rank}(u\varphi)\le1,
 \qquad \|B\|_2<\sqrt{180/\mathfrak c}.
\end{equation}
The line represented by \(u\varphi\) may depend on the witness.
\end{theorem}

\section{Shape-sensitive determinant margins}

The proof of Theorem~\ref{thm:main} uses the full coefficient certificates of
the global-sign theorem, but extracts information not used by the earlier
prime-only estimate.

\subsection{Exterior channel}

For \eqref{eq:Echart}, let \(P_E(u,w)\in\mathbb Z[u,w]\) be defined by
\begin{align}
 -64u^2P_E(u,w)=\mathfrak N\bigl(&16wu,4u(w+1),\notag\\
 &(w+1)(4wu+w+1),4wu(w+1+4wu)\bigr).\label{eq:PE}
\end{align}
The exact division is part of the supplied certificate. With
\[
 J_E=16w^2u^2+8w^2u+28wu+4u+w^2+2w+1,
\]
the literal sum and determinant numerator are
\begin{equation}\label{eq:Escale}
 \frac Sp=\frac{J_E}{16wu},
 \qquad
 -\frac{F}{p^2}=\frac{P_E}{4w^2J_E^6},
 \qquad F=\frac{\mathfrak N}{S^6}.
\end{equation}

The translated polynomial \(P_E(1+X,1+Y)\) has one positive integer
coefficient at every point of the rectangle
\(0\le i\le12,0\le j\le16\), and its least coefficient is \(83886080\).
The translated polynomial
\[
 \bigl[4w^4J_E(u,w)^6\bigr]_{u=1+X,w=1+Y}
\]
has nonnegative coefficients in the same rectangle and coefficient sum
\(64\,449^6\). The verifier reconstructs both objects from \eqref{eq:N} and
checks, coefficient by coefficient, that
\begin{equation}\label{eq:Etranslated}
 449^6P_E(1+X,1+Y)
 -1310720\bigl[4w^4J_E^6\bigr]_{u=1+X,w=1+Y}
\end{equation}
has 221 positive coefficients. Notice that this coefficient statement is in
the translated variables. The unshifted \(P_E\) has negative coefficients;
no unshifted coefficient claim is used.

Since \(u,w\ge1\), \eqref{eq:Etranslated} gives the pointwise inequality
\[
 P_E(u,w)>\frac{1310720}{449^6}\,4w^4J_E(u,w)^6.
\]
The rational coefficient is larger than \(16\mathfrak c_0\). Using
\eqref{eq:Escale} proves
\begin{equation}\label{eq:E-margin}
 -F>\mathfrak c_0p^2w^2>\mathfrak c p^2w^2.
\end{equation}

\subsection{Middle channel}

Set \(b=2+X\), \(c=3+X+Y\), and
\begin{align*}
 L_M&=4X^2+4XY+18X+7Y+19,\\
 J_M&=8X^3+12X^2Y+4XY^2+61X^2+57XY+7Y^2\\
 &\hspace{16mm}+151X+63Y+120.
\end{align*}
The exact middle numerator \(H_M\in\mathbb Z[X,Y]\) is the polynomial defined
by
\[
 -\frac{\mathfrak N}{p^8}=\frac{H_M}{L_M^6},
 \qquad
 \frac Sp=\frac{J_M}{L_M},
 \qquad
 -\frac F{p^2}=\frac{H_M}{J_M^6}.
\]
Its support consists of all 192 lattice points
\[
 0\le i\le18,\qquad0\le j\le12,\qquad i+j\le20,
\]
and every coefficient is positive, with minimum \(81920\). Exact expansion
also proves that every coefficient of
\begin{equation}\label{eq:Mbase-residual}
 483^6H_M-81920J_M^6
\end{equation}
is positive. Therefore
\begin{equation}\label{eq:Mbase}
 -F>\mathfrak c_0p^2.
\end{equation}

Put \(h=Z-Y=p(c-b)\) and \(\tau=h/S\). The product
\((b+c)^2H_M=(5+2X+Y)^2H_M\) has all 237 coefficients in the enlarged
trapezoid
\[
 0\le i\le20,\qquad0\le j\le14,\qquad i+j\le22,
\]
and every coefficient is at least \(81920\). The comparison polynomial
\(L_M^4(Y+1)^2J_M^4\) has nonnegative coefficients, lies in the same
trapezoid, and has coefficient sum \(4\,52^4\,483^4\). The verifier checks
the stronger exact assertion that every coefficient of
\begin{equation}\label{eq:Mgap-residual}
 4\,52^4\,483^4(b+c)^2H_M
 -81920L_M^4(Y+1)^2J_M^4
\end{equation}
is positive. Since
\[
 w=\frac{L_M}{b+c},
 \qquad
 \tau=\frac{(c-b)L_M}{J_M},
\]
this gives
\begin{equation}\label{eq:M-margin}
 -F>\mathfrak c p^2\max\{1,w^2\tau^2\}.
\end{equation}

\section{Cofactors and compound matrices}

Put \(P=\max(p,Y)\) and \(D_*=Z-a\). The reciprocal equation gives
\begin{equation}\label{eq:P-bounds}
 Y\le\frac{2pa}{R}<pw,\qquad p\le P\le pw,\qquad P\le S.
\end{equation}
If \(Z\le2a\), then \(Y\le Z\le2a\) would imply
\(1/a+1/Y+1/Z\ge2/a>4/p\); hence \(Z>2a\) and \(a<D_*\). All four roots
are below \(Z<2D_*\), so \(S<8D_*\). At any root one of the three integer
differences has magnitude at least \(D_*/2\), and the other two have magnitude
at least one. Therefore
\begin{equation}\label{eq:d-xi}
 |d_j|>\frac1{16},\qquad |\xi_j|<4.
\end{equation}
The direct derivative bounds are
\[
 |d_p|,|d_a|\le pP,\qquad |d_Y|\le P^2,\qquad |d_Z|\le S^2.
\]
Substitution in \eqref{eq:receiver} gives the entry table
\begin{equation}\label{eq:entry-table}
\begin{array}{c|ccc}
 &p,a\text{ columns}&Y\text{ column}&Z\text{ column}\\ \hline
0&4&4&4\\
1&\sqrt{pP}&P&S\\
2&3p\sqrt{pP}&3P^2&3S^2\\
3&20pP\sqrt{pP}&20P^3&20S^3.
\end{array}
\end{equation}

In the exterior channel, each \(3\times3\) cofactor has six determinant terms,
each bounded by \(60S^3P^{5/2}\sqrt p\). Thus
\[
 \|\operatorname{adj}O\|_2
 \le\|\operatorname{adj}O\|_F
 \le1440S^3P^{5/2}\sqrt p.
\]
By \eqref{eq:detO} and \eqref{eq:E-margin},
\[
 \|O^{-1}\|_2
 <\frac{360P^{5/2}}{\mathfrak c p^{3/2}w^2}
 \le\frac{360}{\mathfrak c}p\sqrt w.
\]

The middle channel needs its gap and the exact cofactor cancellation. Here
\begin{equation}\label{eq:w-middle}
 w=\frac{4bc}{b+c}-1,\qquad2b-1\le w<4b,\qquad w\ge\frac{19}{5},
\end{equation}
and, with \(d=c-b\ge1\),
\begin{equation}\label{eq:S-h}
 \frac Sh=1+\frac{2b+1+a/p}{d}<2w.
\end{equation}
The three bounds in \eqref{eq:w-middle} retain their exact positive
remainders.  With \(B=b-2\ge0\) and \(D=d-1\ge0\),
\[
 w-(2b-1)=\frac{2bd}{2b+d}>0,
 \qquad
 4b-w=\frac{4b^2+2b+d}{2b+d}>0,
\]
and
\[
 w-\frac{19}{5}
 =\frac{4(5B^2+5BD+13B+4D)}{5(2B+D+5)}\ge0.
\]
After clearing positive denominators, the last strict inequality follows from
\[
 16b^2d+16bd^2-16bd-6d^2-8b^2-6b-3d>0
 \quad(b\ge2,d\ge1).
\]
Indeed, with \(B=b-2\ge0\) and \(D=d-1\ge0\), the left-hand side is exactly
\[
 16B^2D+8B^2+16BD^2+80BD+26B+26D^2+81D+11.
\]
More precisely, writing \(L=4bc-b-c\) and \(\alpha=a/p<1/2\),
\[
 w=\frac{L}{b+c},\qquad
 \frac Sh=1+\frac{2b+1+\alpha}{d}
 <1+\frac{2b+3/2}{d},
\]
and the difference between \(2w\) and the final upper bound is the displayed
positive polynomial divided by \(2d(2b+d)\).
The derivative products give
\[
 |d_a|>\frac{pY}{8},\qquad |d_p|>\frac{pY}{15},\qquad
 |d_Y|\ge\frac{Y^2h}{4S},\qquad
 |d_Z|\ge\frac{Z^2h}{4S}.
\]
For completeness, after division by \(p^2\), the first two left-hand sides are
\[
 \frac{(1-\alpha)(b-\alpha)(c-\alpha)}{1+\alpha+b+c},\qquad
 \frac{(1-\alpha)(b-1)(c-1)}{1+\alpha+b+c}.
\]
The estimates \(b-\alpha>3b/4\), \(b-1\ge b/2\),
\(3(c-\alpha)>1+\alpha+b+c\), and
\(15(c-1)>4(1+\alpha+b+c)\) prove the claimed lower bounds. The tail
bounds follow from
\((b-\alpha)(b-1)>b^2/4\) and
\((c-\alpha)(c-1)>c^2/4\).
For reference, the two denominator comparisons used here are themselves
strict:
\[
 3(c-\alpha)-(1+\alpha+b+c)=2c-b-1-4\alpha>0,
\]
\[
 15(c-1)-4(1+\alpha+b+c)=11c-4b-19-4\alpha>0.
\]
Together with \(b>w/4\), these imply
\begin{equation}\label{eq:xi-middle}
 |\xi_j|\le\frac{12}{p\sqrt w}.
\end{equation}

For three retained columns \(I\), let \(t\) be the omitted root and
\(\Delta_I=\prod_{j<l,\,j,l\in I}(\ell_l-\ell_j)\). Factoring
\(\xi_jd_j=\xi_j^{-1}\) from the three columns, and replacing the final
polynomial row by
\[
 \bigl(7T^2-13H'(T)\bigr)+\frac{13}{A}\bigl(AH'(T)+2T\bigr)
 =7T^2+\frac{26}{A}T,
\]
gives the exact identity
\begin{equation}\label{eq:cofactor}
 \det O_{\{1,2,3\},I}
 =\left(\prod_{j\in I}\xi_j^{-1}\right)\Delta_I
 \left[14AB-64-104A\sum_{j\in I}\ell_j
 -28A^2\sum_{j<l\atop j,l\in I}\ell_j\ell_l\right].
\end{equation}
Using \(A=-1/S\), \(B=-e_2/S\), and the full set of four roots, the bracket
in \eqref{eq:cofactor} becomes
\begin{equation}\label{eq:omitted-bracket}
 40-\frac{76t}{S}-\frac{28t^2}{S^2}-\frac{14e_2}{S^2},
\end{equation}
whose modulus is below \(151\). Moreover,
\begin{equation}\label{eq:xi-product}
 \left|\prod_{j\in I}\xi_j^{-1}\right|
 =S^{-3/2}|\Delta_I|
 \sqrt{\prod_{j\in I}|t-\ell_j|}.
\end{equation}
For an omitted core root, use
\(|\Delta_I|\le PSh\) and
\(\prod_{j\in I}|t-\ell_j|\le pPS\). For an omitted tail, use
\(|\Delta_I|\le pS^2\) and respectively
\(\prod_{j\in I}|t-\ell_j|\le P^2h\) or \(S^2h\).
Equations \eqref{eq:M-margin}, \eqref{eq:cofactor}, and
\eqref{eq:xi-product} bound every entry in the first inverse column by
\begin{equation}\label{eq:first-inverse-column}
 \frac{108}{\mathfrak c}p\sqrt w.
\end{equation}
The constants before \(108\) in the three omitted-root cases are
\(151/4\), \(151\,2^{3/2}/4\), and \(151/(8\sqrt2)\).

For the remaining inverse columns, \eqref{eq:entry-table},
\eqref{eq:xi-middle}, and the baseline part of \eqref{eq:M-margin} give
\[
 \frac{1080w^{3/2}}{\mathfrak c p},\qquad
 \frac{360\sqrt w}{\mathfrak c p^2},\qquad
 \frac{54}{\mathfrak c p^3\sqrt w}.
\]
Since \(p\ge13\) and \(w\le p/3\), these are at most \(28,1,1\) times
\(p\sqrt w/\mathfrak c\). The Frobenius estimate is therefore
\[
 \|O^{-1}\|_2
 <2\sqrt{108^2+28^2+1+1}\,\frac{p\sqrt w}{\mathfrak c}
 <\frac{224}{\mathfrak c}p\sqrt w,
\]
which is covered by \eqref{eq:inverse-main}.

The entries of \(\wedge^2O^{-1}\) are the complementary \(2\times2\)
minors of \(O\), divided by \(\det O\), with their retained signs. In the
exterior channel, every numerator is at most \(120S^3P^2\), hence every
inverse minor is below \(30/\mathfrak c\).

In the middle channel, a tail entry acquires a factor \(\sqrt\tau\) through
\[
 \sqrt{|d_Y|}\le Y\sqrt\tau,\qquad
 \sqrt{|d_Z|}\le Z\sqrt\tau.
\]
Every core--tail inverse minor is bounded by
\begin{equation}\label{eq:core-tail}
 \frac{30\sqrt{w\tau}}
 {\mathfrak c\max\{1,(w\tau)^2\}}
 \le\frac{30}{\mathfrak c},
\end{equation}
and the core--core minors obey the same final bound.

For the tail--tail pair set \(Q_c(T)=(T-a)(T-p)\). Then
\[
 d_Y=\frac{hQ_c(Y)}S,\qquad d_Z=-\frac{hQ_c(Z)}S.
\]
The modulus of the rows-\((2,3)\) minor is exactly
\begin{align}
 \frac{h^2}{S}\sqrt{Q_c(Y)Q_c(Z)}\left|
 14YZ+\frac{26}{S}\{YQ_c(Z)+ZQ_c(Y)\}
 -\frac7{S^2}\{Z^2Q_c(Y)+Y^2Q_c(Z)\}
 \right|.\label{eq:tail-tail}
\end{align}
Using \(Q_c(Y)<Y^2\), \(Q_c(Z)<Z^2\), \(Y+Z<S\), and
\(YZ\le S^2/4\), the bracket is below \(44YZ\); division by
\eqref{eq:M-margin} gives \(11/\mathfrak c\). The row pairs \((1,3)\) and
\((1,2)\) have exact polynomial differences bounded by \(20hS\) and \(3h\),
giving \(5/\mathfrak c\) and \(1/\mathfrak c\). Row-\(0\) pairs are below
\(1/\mathfrak c\) by \eqref{eq:xi-middle}. Thus all 36 entries of the
six-by-six second inverse compound are below \(30/\mathfrak c\), and
\eqref{eq:wedge2-main} follows from its Frobenius norm.

Finally, \eqref{eq:entry-table} gives
\[
 \|O\|_2^2\le\|O\|_F^2
 <64+4S^2+36S^4+1600S^6<41^2S^6.
\]
For a \(4\times4\) matrix,
\(\|\wedge^3O^{-1}\|_2=\|O\|_2/|\det O|\), so
\eqref{eq:wedge3-main} follows from either determinant margin.

It remains to bound \(S\). In E write \(4u+1=tR\). Since
\(R\equiv3\pmod4\), \(t\equiv3\pmod4\) and \(t\ge3\), while direct
substitution gives
\[
 Z=\frac{p(pt+1)}4>p^{3/2}.
\]
In M, \(p\nmid a\) and \(a(4bc-b-c)=pbc\), so
\(p\mid4bc-b-c<4bc\). Hence \(bc>p/4\) and
\[
 S>p(b+c)\ge2p\sqrt{bc}>p^{3/2}.
\]
This proves \eqref{eq:detinv-main}. Since
\[
 \|\wedge^2O^{-1}\|_2=\frac1{s_3(O)s_4(O)}<\frac{180}{\mathfrak c}
\]
and \(s_3\ge s_4\), one has \(s_3>\sqrt{\mathfrak c/180}\).
Removing the largest singular summand of \(O^{-1}\) leaves a rank-three
operator \(B\) with norm \(1/s_3\), proving \eqref{eq:rankone}.

\begin{figure}[t]
\centering
\begin{tikzpicture}[x=1.65cm,y=0.39cm,>=Latex]
  \draw[->] (0.65,-3.4) -- (4.45,-3.4) node[right] {\(\text{singular index }j\)};
  \draw[->] (0.8,-3.8) -- (0.8,13.2) node[above] {\(\alpha_j\) in \(s_j(O_p)\asymp p^{\alpha_j}\)};
  \foreach \j in {1,2,3,4} {
    \draw (\j,-3.55) -- (\j,-3.25);
    \node[below] at (\j,-3.55) {\(\j\)};
  }
  \foreach \y in {-3,0,3,6,9,12} {
    \draw (0.7,\y) -- (0.9,\y);
    \node[left] at (0.7,\y) {\(\y\)};
  }
  \draw[very thick,blue!70!black] plot coordinates {(1,12) (2,4) (3,1.5) (4,-1.5)};
  \foreach \x/\y in {1/12,2/4,3/1.5,4/-1.5}
    \fill[blue!70!black] (\x,\y) circle (2.2pt);
  \draw[very thick,orange!85!black,dashed] plot coordinates {(1,6) (2,2) (3,1) (4,-1)};
  \foreach \x/\y in {1/6,2/2,3/1,4/-1}
    \fill[orange!85!black] (\x,\y) circle (2.2pt);
  \node[anchor=west,blue!70!black] at (1.2,11.1)
    {\(\displaystyle (p+3)/4,\ (pa+2)/3,\ pa(pa+2)/6\)};
  \node[anchor=west,orange!85!black] at (1.2,6.8)
    {\(\displaystyle (3p+1)/8,\ 2a,\ 2pa\)};
  \draw[densely dotted] (0.8,0) -- (4.3,0);
  \node[align=left,anchor=west,fill=white,inner sep=1.5pt] at (1.15,-2.3)
    {\(\alpha_4<0\) is the single shrinking receiver direction;\\
     the two families attain \(p^{3/2}\), \(p^0\), and \(p^{-2}\)
     for the first three inverse compounds.};
\end{tikzpicture}
\caption{Exact singular-value exponents on the two infinite prime families.
The first family attains the \(p^{3/2}\) inverse exponent and the constant
second inverse compound; the second attains the \(p^{-2}\) third inverse
compound exponent. The figure records powers only; the constants are given in
\eqref{eq:family1-singular} and \eqref{eq:family2-singular}.}
\label{fig:singular-exponents}
\end{figure}


\section{Two infinite prime families and exact leading constants}

\subsection{The family through \((13;4,18,468)\)}

For each prime \(p\equiv13\pmod {24}\), put
\begin{equation}\label{eq:family1}
 a=\frac{p+3}{4},\qquad T=pa,\qquad
 Y=\frac{T+2}{3},\qquad Z=\frac{T(T+2)}6.
\end{equation}
Here \(a\) is even and \(T+2\equiv0\pmod3\); hence every displayed quantity
is integral. Moreover
\[
 \frac1Y+\frac1Z=\frac3T,\qquad
 \frac1a+\frac3{pa}=\frac4p.
\]
This is \eqref{eq:Echart} with \(R=3\) and \(u=a^2/2\).

Choose positive real square roots for \(d_a,d_Y>0\) and positive-imaginary
ones for \(d_p,d_Z<0\); then \(\epsilon=1\). In the root order
\((p,a,Y,Z)\), the exact rational root formulas give the following entrywise
asymptotic matrix, with relative error \(O(p^{-1})\) in every displayed
nonzero entry:
\[
O_p=
\begin{pmatrix}
4ip^{-3/2}&4p^{-3/2}&12p^{-2}&96ip^{-4}\\
-\frac14p^{3/2}&-\frac i4p^{3/2}&-\frac i{12}p^2&-\frac1{96}p^4\\
-\frac i2p^{5/2}&\frac18p^{5/2}&\frac1{72}p^4&-\frac i{3072}p^8\\
\frac{13}{64}p^{9/2}&-\frac{13i}{64}p^{9/2}&-\frac i{288}p^6&
\frac5{221184}p^{12}
\end{pmatrix}.
\]
The complete compound matrices give
\begin{equation}\label{eq:family1-compounds}
\begin{array}{c|c|c}
j&\text{power of }p&\text{leading coefficient of }\|\wedge^jO_p\|_2\\ \hline
1&12&5/221184\\
2&16&5/15925248\\
3&35/2&5\sqrt2/(2^{18}3^5)\\
4&16&5/7962624.
\end{array}
\end{equation}
At rank three, two column triples have the same leading order, with
coefficients \(-5/63700992\) and \(-5i/63700992\). Their joint Euclidean norm,
not either selected minor, produces the factor \(\sqrt2\).

Successive compound-norm ratios yield
\begin{equation}\label{eq:family1-singular}
\begin{aligned}
s_1(O_p)&=\frac5{221184}p^{12}(1+O(p^{-1})),&
s_2(O_p)&=\frac1{72}p^4(1+O(p^{-1})),\\
s_3(O_p)&=\frac1{2\sqrt2}p^{3/2}(1+O(p^{-1})),&
s_4(O_p)&=4\sqrt2\,p^{-3/2}(1+O(p^{-1})).
\end{aligned}
\end{equation}
In particular,
\begin{equation}\label{eq:family1-inverse}
 p^{-3/2}\|O_p^{-1}\|_2\longrightarrow\frac1{4\sqrt2},
 \qquad
 \|\wedge^2O_p^{-1}\|_2\longrightarrow\frac12,
 \qquad
 p^4\|\wedge^3O_p^{-1}\|_2\longrightarrow36.
\end{equation}

The leading cofactors give the operator limit
\begin{equation}\label{eq:rankone-limit}
 p^{-3/2}O_p^{-1}
 \longrightarrow\frac18(-i,1,0,0)^Te_0^*.
\end{equation}
Thus the smallest right singular vector tends, up to phase, to
\(v=2^{-1/2}(-i,1,0,0)^T\), and \(p^3s_4(O_p)^2\to32\). The other three
eigenvalues of \(p^3O_p^*O_p\) diverge, so the continuous functional calculus
proves, for every fixed \(t>0\),
\begin{equation}\label{eq:heat-limit}
 e^{-tp^3O_p^*O_p}\longrightarrow e^{-32t}vv^*.
\end{equation}
No statement at \(t=0\) follows from this limit.

\subsection{The family through \((13;5,10,130)\)}

For the same prime progression put
\begin{equation}\label{eq:family2}
 a=\frac{3p+1}{8},\qquad Y=2a,\qquad Z=2pa.
\end{equation}
Then \(a,Y,Z\) are positive integers and
\[
 \frac1a+\frac1{2a}+\frac1{2pa}=\frac4p.
\]
This is \eqref{eq:Echart} with \(R=(p+1)/2\) and \(u=a\), since
\(4u+1=3R\). The derivative signs are
\((+,+,-,-)\); choose positive real and positive-imaginary roots accordingly,
so \(\epsilon=-1\).

The complete leading receiving matrix is
\[
\begin{pmatrix}
\sqrt{32/5}\,p^{-1}&8p^{-1}/\sqrt{15}&i\sqrt{32/3}\,p^{-1}&4ip^{-2}/3\\
-i\sqrt{10}\,p/8&-i\sqrt{15}\,p/8&-\sqrt6\,p/8&-3p^2/4\\
\sqrt{10}\,p^2/4&3\sqrt{15}\,p^2/32&-3i\sqrt6\,p^2/16&-27ip^4/16\\
159i\sqrt{10}\,p^3/256&-33i\sqrt{15}\,p^3/128&
165\sqrt6\,p^3/256&135p^6/16
\end{pmatrix},
\]
again with relative \(O(p^{-1})\) error in each displayed nonzero entry.
The full compound norms are
\begin{equation}\label{eq:family2-compounds}
\begin{array}{c|c|c}
j&\text{power of }p&\text{leading coefficient of }\|\wedge^jO_p\|_2\\ \hline
1&6&135/16\\
2&8&135\sqrt{991}/512\\
3&9&675\sqrt3/512\\
4&8&2565/256.
\end{array}
\end{equation}
Here the relevant small-root row norm, two-row exterior norm, and determinant
are exactly \(991/1024\), \(75/1024\), and \(-19/16\), respectively. Therefore
\begin{equation}\label{eq:family2-singular}
\begin{aligned}
s_1(O_p)&\sim\frac{135}{16}p^6,&
s_2(O_p)&\sim\frac{\sqrt{991}}{32}p^2,\\
s_3(O_p)&\sim\frac{5\sqrt3}{\sqrt{991}}p,&
s_4(O_p)&\sim\frac{38}{5\sqrt3}p^{-1},
\end{aligned}
\end{equation}
and
\begin{equation}\label{eq:family2-inverse}
 p^{-1}\|O_p^{-1}\|_2\to\frac{5\sqrt3}{38},
 \qquad
 \|\wedge^2O_p^{-1}\|_2\to\frac{\sqrt{991}}{38},
 \qquad
 p^2\|\wedge^3O_p^{-1}\|_2\to\frac{16}{19}.
\end{equation}

Dirichlet's theorem applies because \(\gcd(13,24)=1\), so both families
contain infinitely many prime witnesses. Equations
\eqref{eq:family1-inverse} and \eqref{eq:family2-inverse} prove the optimality
of all three exponents in Theorem~\ref{thm:main}.

\section{Metrics, conductor, and observations}\label{sec:transport}

Let \(G\) and \(Q\) be positive-definite Hermitian matrices on
\(V_\beta=\mathbb C^4_{\mathrm{labels}}\) and
\(E_q=\mathbb C^4_q\), respectively. For \(1\le j\le4\), exterior-power
submultiplicativity applied to
\(G^{1/2}O^{-1}Q^{-1/2}\) gives
\begin{equation}\label{eq:metric-wedge}
 \|\wedge^jO^{-1}\|_{Q\to G}
 \le
 \left[
 \frac{\prod_{a=1}^j\lambda_a^\downarrow(G)}
 {\prod_{a=1}^j\lambda_a^\uparrow(Q)}
 \right]^{1/2}B_j(p),
\end{equation}
where \(B_j\) is the corresponding bound in
\eqref{eq:inverse-main}--\eqref{eq:detinv-main}. If the metrics depend on
\(p\), their eigenvalue factors remain present.

The conductor maps in the source programme have the exact types
\[
 O:V_\beta\to E_q,\qquad
 \Phi:E_q\to U,\qquad
 L_U:U\to W,\qquad
 \Psi=L_U\Phi:E_q\to W.
\]
On the stated invertible restrictions, the complete composite and inverse are
\begin{equation}\label{eq:conductor-composite}
 C=L_U\Phi O=\Psi O,
 \qquad
 C^{-1}=O^{-1}\Phi^{-1}L_U^{-1}=O^{-1}\Psi^{-1}.
\end{equation}
Thus \eqref{eq:rankone} transports as
\begin{equation}\label{eq:rankone-transport}
 C^{-1}=B\Psi^{-1}+u(\varphi\Psi^{-1}).
\end{equation}
The second summand has rank at most one. The first is uniformly bounded only
when \(\Psi^{-1}\) is fixed or uniformly bounded; the receiver estimate does
not remove conductor resonances.

For a further linear observation \(\Lambda:E_q\to Y\), let
\(Y_0=\operatorname{im}(\Lambda O)=\operatorname{im}\Lambda\). Then
\begin{equation}\label{eq:observation}
 \ker(\Lambda O)=O^{-1}\ker\Lambda,
 \qquad
 \operatorname{rank}(\Lambda O)=\operatorname{rank}\Lambda,
\end{equation}
and the quotient covariance on the actual image is
\begin{equation}\label{eq:covariance}
 Q_{\Lambda O}
 =\left[(\Lambda OG^{-1}O^*\Lambda^*)|_{Y_0}\right]^{-1}.
\end{equation}
In an ambient codomain larger than \(Y_0\), the corresponding formula uses
the Moore--Penrose inverse. An observation after the conductor retains every
factor:
\[
 \ker(\Lambda C)=O^{-1}\Phi^{-1}L_U^{-1}\ker\Lambda.
\]
The metric paired exactly to \(O\) is
\[
 G_{\mathrm{out}}=O^{-*}G_{\mathrm{in}}O^{-1},
\]
and through the conductor it is
\[
 C^{-*}G_{\mathrm{in}}C^{-1}
 =\Psi^{-*}O^{-*}G_{\mathrm{in}}O^{-1}\Psi^{-1}.
\]
Neither \(\Phi\) nor \(\Psi\) is declared unitary. Finally,
\(\epsilon\det O=4(-F)S^3>0\), but no sign or phase for
\(\det(\Psi O)\) follows without retaining \(\det\Psi\).

\begin{figure}[t]
\centering
\resizebox{\textwidth}{!}{%
\begin{tikzpicture}[
  node distance=15mm and 18mm,
  box/.style={draw,rounded corners,minimum width=30mm,minimum height=10mm,align=center},
  >=Latex
]
  \node[box] (V) {\(V_\beta=\mathbb C^4_{\rm labels}\)};
  \node[box,right=of V] (E) {\(E_q=\mathbb C^4_q\)};
  \node[box,right=of E] (U) {\(U=\mathcal P_{v+3}/\mathcal P_{v-1}\)};
  \node[box,right=of U] (W) {\(W=\mathcal P_3\)};
  \draw[->,thick] (V) -- node[above] {\(O\)} (E);
  \draw[->,thick] (E) -- node[above] {\(\Phi\)} (U);
  \draw[->,thick] (U) -- node[above] {\(L_U\)} (W);
  \draw[->,thick,bend right=23] (E) to node[below] {\(\Psi=L_U\Phi\)} (W);
  \draw[->,very thick,bend left=24] (V) to
    node[above] {\(C=L_U\Phi O=\Psi O\)} (W);
  \node[box,below=of E] (Y) {\(Y_0=\operatorname{im}(\Lambda O)\)};
  \draw[->,thick] (E) -- node[right] {\(\Lambda|_{\operatorname{im}O}\)} (Y);
  \node[align=left,below=of U,draw,dashed,rounded corners,inner sep=4pt] (inv)
    {\(C^{-1}=O^{-1}\Phi^{-1}L_U^{-1}=O^{-1}\Psi^{-1}\)\\
     \(\ker(\Lambda C)=O^{-1}\Phi^{-1}L_U^{-1}\ker\Lambda\)};
  \draw[densely dotted] (W) -- (inv);
\end{tikzpicture}%
}
\caption{The exact typed conductor chain. Omitting \(\Phi\) makes the composite
ill-typed. The receiver's rank-one inverse decomposition passes through
\(\Psi^{-1}\), while any growth or resonance of \(\Psi^{-1}\) remains.}
\label{fig:typed-transport}
\end{figure}


\section{Scope}

The theorem begins with an existing integral witness. It proves that the
literal receiving stage is invertible there, bounds its inverse and compound
maps, and determines sharp growth on two infinite families. It does not
produce a witness at a prescribed hard prime and therefore does not prove the
Erd\H{o}s--Straus conjecture. The heat limit
\eqref{eq:heat-limit} belongs to the displayed finite-dimensional ES receiver;
it is not an RH heat operator. The conductor conclusions are exactly those
of \eqref{eq:conductor-composite}--\eqref{eq:rankone-transport}, with all
metric and resonance factors retained.
